\chapter{Prerequisites}
\label{ch:prereq}

Before building Veilbox, the host system must meet specific hardware and
software requirements. This chapter documents every prerequisite, from
the Fedora version and kernel release to the precise toolchain versions
needed for successful compilation.

\section{Hardware Requirements}

\subsection{Build Machine}

Building the Veilbox kernel and initramfs requires a machine with the
following minimum specifications:

\begin{itemize}[nosep]
  \item \textbf{CPU}: x86\_64, dual-core or better. Kernel compilation is
        heavily parallelized (make -j\$(nproc)), so more cores dramatically
        reduce build time. A four-core CPU completes a full build in
        approximately 5 minutes; a single-core CPU may take 20+ minutes.
  \item \textbf{RAM}: 4 GB minimum, 8 GB recommended. The kernel build
        process spawns multiple GCC instances, each consuming 200--500 MB
        of memory. The embedded initramfs (approximately 90 MB when
        uncompressed) must fit in the kernel's temporary build directory.
  \item \textbf{Storage}: 20 GB free. The Linux kernel source tree occupies
        approximately 6 GB after extraction (or 3.5 GB as a Git clone with
        history). Build artifacts (object files, .o files, final binaries)
        add another 4--6 GB. The output disk image is 8 GB.
  \item \textbf{Network}: Broadband internet connection for downloading
        kernel source, BusyBox, containerd, runc, nerdctl, and Dropbear
        tarballs. Total download size is approximately 200 MB.
\end{itemize}

\subsection{Test/Runtime Machine}

The Veilbox disk image runs on any x86\_64 machine with:

\begin{itemize}[nosep]
  \item \textbf{RAM}: 2 GB minimum. The initramfs occupies approximately
        90 MB of RAM; the kernel uses about 50 MB; containerd, runc, and
        Dropbear add another 50 MB combined. The remaining memory is
        available for containers.
  \item \textbf{Storage}: 8 GB disk (raw or virtual). The VEILBOX partition
        uses approximately 500 MB for the kernel, GRUB modules, and state
        directory. The remaining ~6.5 GB is unused and available for
        container images and volumes on the state partition.
  \item \textbf{Display}: VGA text console (1024x768 minimum) or serial
        console (9600 baud, 8N1). No graphical output is required.
  \item \textbf{Network}: Any Ethernet adapter (virtio-net, e1000, rtl8139
        for VMs; Intel e1000e, Realtek for bare metal).
\end{itemize}

\section{Software Requirements}

\subsection{Host Operating System}

The build script was developed and tested on Fedora 44, but any Linux
distribution with GCC 14+, GNU Make 4.4+, and Python 3.10+ should work.

\begin{longtable}{|p{4cm}|p{4cm}|p{4cm}|}
\hline
\textbf{Distribution} & \textbf{Status} & \textbf{Notes} \\
\hline
Fedora 44 & Verified & Primary development platform. \\
Fedora 43 & Compatible & May require updated packages. \\
RHEL 9 / CentOS 9 & Compatible & EPEL repository needed. \\
Debian 12 / Ubuntu 24.04 & Compatible & Package names differ. \\
Arch Linux & Compatible & Use distro-specific package names. \\
\hline
\end{longtable}

\subsection{Kernel Source}

The kernel source is obtained from kernel.org or a Git mirror:

\begin{itemize}[nosep]
  \item \textbf{Base version}: 7.1.0-rc6 (verified working).
  \item \textbf{Source URL}: \url{https://cdn.kernel.org/pub/linux/kernel/v7.x/linux-7.1.0-rc6.tar.xz}
  \item \textbf{Size}: Approximately 145 MB compressed, 1.1 GB extracted.
  \item \textbf{Alternate source}: Git clone from \url{https://git.kernel.org/pub/scm/linux/kernel/git/torvalds/linux.git}
\end{itemize}

The build script (\texttt{build.sh}) downloads and extracts the kernel
automatically. It checks for an existing \texttt{linux-7.1.0-rc6/} directory
and skips extraction if present.

\subsection{Toolchain}

The following tools must be installed on the build machine:

\begin{longtable}{|p{4cm}|p{3cm}|p{5cm}|}
\hline
\textbf{Tool} & \textbf{Version} & \textbf{Purpose} \\
\hline
GCC & 14+ & Kernel and BusyBox compilation. \\
GNU Make & 4.4+ & Build system orchestration. \\
GNU Binutils & 2.46+ & Assembly, linking, object manipulation. \\
flex & 2.6+ & Kernel configuration parser. \\
bison & 3.8+ & Kernel configuration parser. \\
bc & 1.07+ & Kernel build-time calculations. \\
OpenSSL dev & 3.0+ & Kernel module signing headers. \\
elfutils-libelf-devel & 0.190+ & Kernel build headers (libelf). \\
perl & 5.36+ & Kernel build scripts. \\
python3 & 3.10+ & Kernel configuration tools. \\
rsync & 3.2+ & File synchronization. \\
cpio & 2.13+ & Initramfs cpio archive. \\
xz & 5.4+ & Kernel compression support. \\
qemu-img & 7.0+ & Disk image format conversion. \\
debugfs & 1.47+ & Rootless ext4 filesystem population. \\
sfdisk & 2.39+ & Disk partitioning. \\
grub2-tools & 2.06+ & GRUB core.img generation. \\
qemu-system-x86\_64 & 7.0+ & VM testing. \\
\hline
\end{longtable}

\subsection{Installation Commands}

On Fedora:
\begin{lstlisting}
sudo dnf install gcc make flex bison bc openssl-devel \
  elfutils-libelf-devel perl python3 rsync cpio xz \
  qemu-img e2fsprogs util-linux grub2-tools \
  qemu-system-x86-x86_64
\end{lstlisting}

On Debian/Ubuntu:
\begin{lstlisting}
sudo apt install gcc make flex bison bc libssl-dev \
  libelf-dev perl python3 rsync cpio xz-utils \
  qemu-utils e2fsprogs util-linux grub2-common \
  qemu-system-x86
\end{lstlisting}

\subsection{Kernel Configuration Dependencies}

The kernel build system requires several host utilities that are not
always installed by default. The build script's \texttt{install\_deps()}
function checks for each of these and reports missing dependencies before
starting the build.

\begin{itemize}[nosep]
  \item \texttt{hostname}: Used by kernel makefiles for build identification.
  \item \texttt{getconf}: Determines number of CPU cores for parallel builds.
  \item \texttt{tee}: Logging build output to stdout and log file.
  \item \texttt{curl} or \texttt{wget}: Downloading source tarballs.
  \item \texttt{sha256sum}: Verifying downloaded file integrity.
  \item \texttt{bash 5+}: The build script requires bash-specific features
        (arrays, process substitution, \texttt{[[ ]]}) tests.
\end{itemize}

\section{Network Requirements}

\subsection{Build-Time Network Access}

The build script requires internet access for:

\begin{enumerate}[nosep]
  \item Downloading the Linux kernel source tarball.
  \item Downloading BusyBox source from busybox.net.
  \item Downloading containerd, runc, and nerdctl from GitHub releases.
  \item Downloading Dropbear SSH from dropbear.nl.
  \item (Optional) Cloning kernel Git repository.
\end{enumerate}

All downloads are cached in the host \$HOME/veilbox-cache directory
for offline builds.

\subsection{Runtime Network (Guest)}

The guest OS requires:
\begin{itemize}[nosep]
  \item A DHCP server on the virtual network. QEMU provides this via
        SLiRP (user-mode networking) by default.
  \item Internet access via NAT or bridged networking for container image
        pulls.
  \item Port 22 reachable on the guest IP for SSH access.
\end{itemize}

\section{Virtualization Requirements}

\subsection{QEMU}

QEMU version 7.0 or later with KVM support is recommended for testing.
The test.sh script expects:

\begin{itemize}[nosep]
  \item \texttt{qemu-system-x86\_64} binary in PATH.
  \item KVM device (\texttt{/dev/kvm}) accessible. On most distributions,
        the user must be in the \texttt{kvm} group.
  \item Approximately 2 GB of available host RAM for the guest VM.
  \item For \texttt{-nographic} mode: a terminal emulator that supports
        the serial console (any modern terminal works).
\end{itemize}

\subsection{VirtualBox (Alternative)}

For VirtualBox deployment:
\begin{itemize}[nosep]
  \item VirtualBox 6.1+ or 7.0+ with Extension Pack (for VirtIO support).
  \item Import the VDI file via File $\rightarrow$ Import Appliance.
  \item Configure: 2 GB RAM, 1 CPU, VirtIO-Net adapter, enable serial
        port (optional).
\end{itemize}

\section{User Privileges}

One of the design goals of Veilbox is a completely rootless build. No
sudo access is required for any stage of the build process:

\begin{itemize}[nosep]
  \item \textbf{Kernel compilation}: make, gcc, and ld all run as the
        regular user.
  \item \textbf{Filesystem creation}: debugfs writes files into the ext4
        image without requiring loop device access (which would need root).
  \item \textbf{GRUB installation}: A Python script writes boot sectors
        directly to the raw disk image file. No loop device mount needed.
  \item \textbf{Disk partitioning}: sfdisk writes the MBR partition table
        to a regular file.
  \item \textbf{Format conversion}: qemu-img reads the raw image and
        creates VDI output.
\end{itemize}

The only exception is dependency installation: \texttt{install\_deps()}
uses sudo to run \texttt{dnf install}. If sudo is not available or
passwordless, the script prints the required package list and continues.
